%% kapa: exam=1 round=20 year=2009 period=2 form=A subject=accounting subject_ko=회계학 seq=20 number=20
%% source: Q-Net 감정평가사 기출문제 (cst00309) artlSeq=5024178 file=extracted/1008/5024178/2009년 제20회 감정평가사 1차시험 기출문제/2교시.hwp
%% conversion: kordoc 4.24.51 -> tools/gen1.py (auto); review=proofread-sonnet; flags=none
\begin{question}{20}
(주)한국은 (주)서울을 흡수합병 하였으며 이 합병은 매수에 해당한다. 합병당시 (주)서울의 식별 가능한 자산과 부채 및 자본의 장부금액과 공정가치는 다음과 같다.

{\small\centering
\begin{tabular}{l r r}
\hline
 & \multicolumn{1}{c}{장부금액} & \multicolumn{1}{c}{공정가치} \\
\hline
자산 & ￦8,500,000 & ￦9,000,000 \\
부채 & 5,000,000 & 4,000,000 \\
\makecell[l]{자본금\\(1주당 액면금액 ￦1,000)} & 3,000,000 & \\
이익잉여금 & 500,000 & \\
\hline
\end{tabular}\par}
(주)한국은 (주)서울을 합병하면서 (주)서울의 주식 1주에 대하여 (주)한국의 주식 1.5주를 발행하여 교부하였다. 합병당시 (주)한국 주식의 1주당 액면가액은 ￦1,000이며, 시장가격은 ￦1,200이다. 이 합병으로 인해 (주)한국이 계상하게 될 영업권(또는 부의영업권)은 얼마인가? \par
\begin{choices}
  \item 영업권 ￦200,000
  \item 영업권 ￦400,000
  \item 부의영업권 ￦500,000
  \item 부의영업권 ￦2,000,000
  \item 부의영업권 ￦2,400,000
\end{choices}
\end{question}
\answer{2}
